%----------------------------------------------------------------------------------------
% KANS. Alle opgaven zijn eigen werk.
%
% Opbouw: Start, leerstof (theorie, werkwijze of uitleg, voorbeeld, korte oefening), en op een
% nieuwe pagina de Opdrachten (door elkaar, stijgende moeilijkheid), Koppel en Test jezelf.
% Labels van dit hoofdstuk beginnen met h1-.
%----------------------------------------------------------------------------------------

\mychapter{1}{Kans}

\lkr{Leerplan B+S'': LPD 48 (kansen met kruistabellen, boomdiagrammen en de wet van Laplace), met het verband tussen relatieve frequentie en kans en de terminologie uitkomst, gebeurtenis en uitkomstenverzameling. Dit hoofdstuk legt de basis: Laplace, de wet van de grote aantallen en de rekenregels. Kruistabellen en boomdiagrammen volgen in hoofdstuk 2. Telproblemen (zelfde jaar) helpt bij het tellen van gunstige uitkomsten; waar die cursus nodig is, staat dat erbij.}

Gooi je met een dobbelsteen, dan weet je niet welk getal er komt. Toch is er orde in het toeval: gooi je heel vaak, dan komt elk getal ongeveer even vaak. Kansrekenen maakt die orde precies. Je gebruikt het om een spel te ontwerpen, een risico te schatten of een medische test te beoordelen. In dit hoofdstuk leer je wat een kans is, hoe je ze berekent als alle uitkomsten even waarschijnlijk zijn, en welke rekenregels altijd gelden.

\begin{oefeningen}[Start]
\oef{h1-start-1}

\oef{h1-verken}
\end{oefeningen}

\lkr{Laat oefening \ref{oef:h1-verken} maken vóór de uitleg. De fout in vraag d is de belangrijkste misvatting van het hoofdstuk: de wet van Laplace geldt alleen als alle uitkomsten even waarschijnlijk zijn, en de 11 sommen zijn dat niet.}

%========================================================================================
\section{Uitkomsten en gebeurtenissen}

\subsection{Uitkomstenverzameling en gebeurtenis}

Een \textbf{kansexperiment} is een experiment waarvan je de uitkomst niet op voorhand kent, zoals een worp met een dobbelsteen of een trekking uit een vaas. Om erover te rekenen, schrijf je eerst op wat er allemaal kan gebeuren. Daarna beschrijf je wat je wilt weten, zoals \textit{een even getal gooien}, als een deel daarvan.

\begin{theorie}
De \textbf{\textcolor{accent}{uitkomstenverzameling}} $\formule{\Omega}$ van een kansexperiment is de verzameling van alle mogelijke uitkomsten. Een \textbf{\textcolor{accent}{gebeurtenis}} is een deelverzameling $\formule{A \subseteq \Omega}$.
\end{theorie}

Een uitkomst is een element van $\Omega$, een gebeurtenis een verzameling van uitkomsten. Bij een worp met een dobbelsteen is $\Omega = \{1, 2, 3, 4, 5, 6\}$, en \textit{een even getal gooien} is de gebeurtenis $\{2, 4, 6\}$. Een gebeurtenis \textit{treedt op} als de uitkomst erin ligt. Twee gebeurtenissen zijn bijzonder: $\Omega$ zelf treedt altijd op (de \textit{zekere} gebeurtenis), en de lege verzameling $\emptyset$ nooit (de \textit{onmogelijke} gebeurtenis, zoals \textit{een zeven gooien}).

\needspace{16\baselineskip}
\subsection{Bewerkingen met gebeurtenissen}

In gewone taal combineer je gebeurtenissen met de woorden \textit{niet}, \textit{en} en \textit{of}: \textit{geen zes}, \textit{even en groter dan 3}. In de taal van verzamelingen zijn dat het complement, de doorsnede en de unie.

\begin{theorie}
{\centering\arrayrulecolor{accent}\renewcommand{\arraystretch}{1.45}%
\begin{tabular}{|>{\centering\arraybackslash}m{3.6cm}|>{\centering\arraybackslash}m{2.6cm}|>{\centering\arraybackslash}m{6.4cm}|}
\hline
\textbf{gebeurtenis} & \textbf{notatie} & \textbf{betekenis} \\ \hline
niet $A$ (complement) & $\formule{\compl{A}}$ & alle uitkomsten die niet in $A$ zitten \\ \hline
$A$ en $B$ & $\formule{A \cap B}$ & de uitkomsten die in $A$ én in $B$ zitten \\ \hline
$A$ of $B$ & $\formule{A \cup B}$ & de uitkomsten in $A$, in $B$, of in beide \\ \hline
\end{tabular}\par}
\end{theorie}

In een venndiagram is $\Omega$ een rechthoek en is elke gebeurtenis een gebied erin. De figuren tonen de drie bewerkingen in het donkere deel.

\figuren{\begin{tikzpicture}[font=\small,scale=.8]
  \fill[accent!40] (-1.3,-1.1) rectangle (2.5,1.1); \fill[white] (0,0) circle (.85);
  \draw[grijs] (-1.3,-1.1) rectangle (2.5,1.1); \draw[accent,thick] (0,0) circle (.85);
  \node[text=accent] at (0,0) {$A$}; \node[below] at (.6,-1.1) {$\compl{A}$};
\end{tikzpicture}\qquad
\begin{tikzpicture}[font=\small,scale=.8]
  \begin{scope}\clip (0,0) circle (.85); \fill[accent!40] (1.15,0) circle (.85); \end{scope}
  \draw[grijs] (-1.3,-1.1) rectangle (2.5,1.1);
  \draw[accent,thick] (0,0) circle (.85); \draw[accent,thick] (1.15,0) circle (.85);
  \node[text=accent] at (-.35,0) {$A$}; \node[text=accent] at (1.5,0) {$B$}; \node[below] at (.6,-1.1) {$A \cap B$};
\end{tikzpicture}\qquad
\begin{tikzpicture}[font=\small,scale=.8]
  \fill[accent!40] (0,0) circle (.85); \fill[accent!40] (1.15,0) circle (.85);
  \draw[grijs] (-1.3,-1.1) rectangle (2.5,1.1);
  \draw[accent,thick] (0,0) circle (.85); \draw[accent,thick] (1.15,0) circle (.85);
  \node[text=accent] at (-.35,0) {$A$}; \node[text=accent] at (1.5,0) {$B$}; \node[below] at (.6,-1.1) {$A \cup B$};
\end{tikzpicture}}

Het woord \textit{of} betekent in de wiskunde altijd \textit{of beide}: de overlap hoort bij $A \cup B$. In de omgangstaal is dat niet altijd zo. \textit{Koffie of thee?} bedoelt meestal één van de twee.

\voorbeeld{}{Je gooit met een dobbelsteen. Noem $A$ \textit{een even getal} en $B$ \textit{minstens 4}. Bepaal $\compl{A}$, $A \cap B$ en $A \cup B$.\label{vb:h1-gebeurtenis}}{%
\uitleg{We schrijven elke gebeurtenis als verzameling van uitkomsten, en lezen de bewerkingen af.}
\naast{\begin{minipage}{10.2cm}
\begin{stappenvrij}{5.2cm}{4.7cm}
\Omega = \{1, 2, 3, 4, 5, 6\} & De uitkomstenverzameling. \\[.6ex]
A = \{2, 4, 6\}, \ B = \{4, 5, 6\} & De twee gebeurtenissen. \\[.6ex]
\compl{A} = \{1, 3, 5\} & Een oneven getal. \\[.6ex]
A \cap B = \{4, 6\} & Even én minstens 4. \\[.6ex]
A \cup B = \{2, 4, 5, 6\} & Even, of minstens 4, of beide.
\end{stappenvrij}
\end{minipage}}{\begin{tikzpicture}[font=\small]
  \draw[grijs] (-1.3,-1.1) rectangle (2.5,1.1); \node[grijs,anchor=north east] at (2.5,1.1) {$\Omega$};
  \draw[verz] (0,0) circle (.85); \draw[verz] (1.15,0) circle (.85);
  \node at (-.4,0) {2}; \node at (.58,.25) {4}; \node at (.58,-.25) {6}; \node at (1.55,0) {5};
  \node at (-1.05,-.85) {1}; \node at (2.25,-.85) {3};
  \node[text=accent] at (-.5,.62) {$A$}; \node[text=accent] at (1.65,.62) {$B$};
\end{tikzpicture}}}

\begin{oefeningen}
\oef{h1-woorden}
\end{oefeningen}

\needspace{16\baselineskip}
\subsection{Onverenigbare gebeurtenissen}

Soms kunnen twee gebeurtenissen niet samen optreden: een worp is nooit tegelijk een 1 en een 6. Dat geval krijgt een eigen naam, want straks rekenen we er het eenvoudigst mee.

\begin{theorie}
$A$ en $B$ zijn \textbf{\textcolor{accent}{onverenigbaar}} (of \textit{disjunct}) als \quad $\formule{A \cap B = \emptyset}$
\end{theorie}

\naast[3.2cm]{In een venndiagram overlappen onverenigbare gebeurtenissen niet. Een gebeurtenis en haar complement zijn altijd onverenigbaar, want een uitkomst ligt nooit tegelijk in $A$ en buiten $A$. Let op: onverenigbaar zegt iets over één uitvoering van het experiment. Gooi je twee keer, dan kan de eerste worp een 1 zijn en de tweede een 6.}{\begin{tikzpicture}[font=\small,scale=.8]
  \draw[grijs] (-1.1,-1) rectangle (2.9,1); \node[grijs,anchor=north east] at (2.9,1) {$\Omega$};
  \draw[verz] (0,0) circle (.7); \draw[verz] (1.8,0) circle (.7);
  \node[text=accent] at (0,0) {$A$}; \node[text=accent] at (1.8,0) {$B$};
\end{tikzpicture}}

\voorbeeld{}{Je trekt een kaart uit een spel van 52. Noem $H$ \textit{een harten}, $S$ \textit{een schoppen} en $A$ \textit{een aas}. Welke van deze gebeurtenissen zijn twee aan twee onverenigbaar?\label{vb:h1-onverenigbaar}}{%
\uitleg{We zoeken voor elk paar een kaart die in beide gebeurtenissen ligt. Bestaat die niet, dan is de doorsnede leeg.}
\begin{stappenbreed}
H \cap S = \emptyset & Een kaart heeft maar één kleur: onverenigbaar. \\[.6ex]
H \cap A = \{\text{hartenaas}\} & Niet leeg: niet onverenigbaar. \\[.6ex]
S \cap A = \{\text{schoppenaas}\} & Ook niet onverenigbaar.
\end{stappenbreed}
Alleen $H$ en $S$ zijn onverenigbaar.}

\begin{oefeningen}
\oef{h1-gebeurtenis}
\end{oefeningen}

%========================================================================================
\section{De wet van Laplace}

\subsection{Even waarschijnlijke uitkomsten}

Bij een eerlijke dobbelsteen heeft elk van de zes getallen dezelfde kans. Samen moeten die kansen 1 zijn (het is zeker dat er een getal komt), dus elk getal heeft kans $\frac16$. Een even getal komt bij drie van die zes uitkomsten, dus de kans erop is $3 \cdot \frac16 = \frac36$. Zo rekende de Franse wiskundige Pierre-Simon Laplace.

Algemeen: zijn er $\#\Omega$ even waarschijnlijke uitkomsten, dan heeft elke uitkomst kans $\frac{1}{\#\Omega}$, en een gebeurtenis met $\#A$ uitkomsten heeft $\#A$ keer die kans.

\begin{theorie}
\begin{stelling}{Wet van Laplace}
Heeft een kansexperiment een eindig aantal uitkomsten die allemaal even waarschijnlijk zijn, dan is de kans op een gebeurtenis $A$

\vspace{1mm}
\qquad $\formule{P(A) = \dfrac{\#A}{\#\Omega} = \dfrac{\text{aantal gunstige uitkomsten}}{\text{aantal mogelijke uitkomsten}}}$
\end{stelling}
\end{theorie}

Een bewijs in de strikte zin is er niet: de wet van Laplace is de afspraak waarmee we kansen toekennen als alle uitkomsten even waarschijnlijk zijn. De redenering hierboven toont waarom het de enige redelijke afspraak is. De uitkomsten in $A$ heten de \textit{gunstige} uitkomsten, ook als $A$ iets vervelends is.

De voorwaarde \textit{even waarschijnlijk} is belangrijk. Kies je de uitkomsten verkeerd, dan geeft de formule een fout antwoord.

\voorbeeld{}{Je gooit twee munten. Wat is de kans op één keer kop en één keer munt?\label{vb:h1-munten}}{%
\uitleg{Een snelle redenering zegt: er zijn drie mogelijkheden (0, 1 of 2 keer kop), dus de kans is $\frac13$. Maar zijn die drie even waarschijnlijk? Maak de munten verschillend, een eerste en een tweede munt.}
\begin{stappenbreed}
\Omega = \{KK,\, KM,\, MK,\, MM\} & Vier rijtjes, en die zijn wel even waarschijnlijk: elke munt valt even makkelijk op kop als op munt. \\[.6ex]
P(\text{één kop}) = \frac{\#\{KM,\, MK\}}{4} = \frac24 = \frac12 & \textit{Eén keer kop} ontstaat op twee manieren, \textit{twee keer kop} maar op één.
\end{stappenbreed}
Het antwoord $\frac13$ is fout: de drie aantallen zijn niet even waarschijnlijk. Gooi je het echt vaak, dan zie je in ongeveer de helft van de worpen één kop.}

Met twee dobbelstenen is het net zo: er zijn 11 mogelijke sommen, van 2 tot 12, maar die zijn niet even waarschijnlijk. De 36 koppels $(a, b)$ zijn dat wel: daarom tel je koppels, met een \textbf{roosterdiagram}.

\voorbeeld{}{Je gooit met twee dobbelstenen. Wat is de kans dat de som 7 is? En dat de som minstens 10 is?\label{vb:h1-laplace}}{%
\uitleg{De 36 koppels zijn even waarschijnlijk. In het rooster lezen we af hoeveel koppels gunstig zijn.}
\naast{\begin{minipage}{9.4cm}
\begin{stappenvrij}{4.2cm}{4.8cm}
\#\Omega = 6 \cdot 6 = 36 & Koppels (eerste, tweede dobbelsteen). \\[1ex]
P(\text{som } 7) = \frac{6}{36} = \frac16 & De diagonaal $(1, 6), (2, 5), \ldots, (6, 1)$. \\[3ex]
P(\text{som} \geqslant 10) = \frac{6}{36} = \frac16 & Som 10: 3 koppels, 11: 2, 12: 1.
\end{stappenvrij}
\end{minipage}}{\begin{tikzpicture}[x=.48cm,y=.48cm,font=\scriptsize]
  \foreach \a in {1,...,6}{\foreach \b in {1,...,6}{
    \pgfmathtruncatemacro{\s}{\a+\b}
    \ifnum\s=7 \fill[accent!35] (\a-.5,\b-.5) rectangle (\a+.5,\b+.5);\fi
    \ifnum\s>9 \fill[fel!45] (\a-.5,\b-.5) rectangle (\a+.5,\b+.5);\fi
    \node at (\a,\b) {\s};}}
  \draw[grijs] (.5,.5) grid (6.5,6.5);
  \foreach \t in {1,...,6}{\node[grijs] at (\t,0) {\t}; \node[grijs] at (0,\t) {\t};}
\end{tikzpicture}}}

\begin{werkwijze}
Zoek een uitkomstenverzameling met even waarschijnlijke uitkomsten. Maak gelijke voorwerpen verschillend (een eerste en een tweede munt) en werk met koppels of rijtjes, niet met sommen of aantallen. Tel de mogelijke en de gunstige uitkomsten, en deel.
\end{werkwijze}

\nujij{Wat is de kans dat het product van de ogen van twee dobbelstenen even is?}{Het product is oneven als beide oneven zijn: $3 \cdot 3 = 9$ koppels. Even: $36 - 9 = 27$, dus $\frac{27}{36} = \frac34$.}

\begin{oefeningen}
\oef{h1-laplace-fout}

\oef{h1-rooster}

\oef{h1-laplace}
\end{oefeningen}

\needspace{16\baselineskip}
\subsection{Tellen met combinaties}

Trek je 3 kaarten uit een spel van 52, dan is een rooster veel te groot. Neem je $k$ elementen tegelijk uit een groep van $n$, zonder op de volgorde te letten, dan is het aantal mogelijke groepen $\binom{n}{k}$, op het rekentoestel \texttt{nCr}. Bij een blinde trekking is elke groep even waarschijnlijk, dus de wet van Laplace geldt.

\voorbeeld{}{Je trekt 3 kaarten uit een spel van 52. Wat is de kans dat het drie harten zijn?\label{vb:h1-kaarten}}{%
\uitleg{Elke groep van 3 kaarten is even waarschijnlijk. De volgorde telt niet, dus tellen we groepen met $\binom{n}{k}$.}
\begin{stappen}
\#\Omega = \binom{52}{3} = 22\,100 & Alle groepen van 3 kaarten uit 52. \\[2ex]
\#A = \binom{13}{3} = 286 & De gunstige groepen: 3 kaarten uit de 13 harten. \\[2ex]
P(A) = \frac{286}{22\,100} = \frac{11}{850} \approx 0{,}013 & Iets meer dan 1~\%.
\end{stappen}}

\begin{werkwijze}
Neem je $k$ elementen tegelijk uit $n$, zonder volgorde: tel alle groepen met $\binom{n}{k}$. Tel de gunstige groepen op dezelfde manier, maar kies alleen uit het gunstige deel van de $n$ elementen. Deel de twee aantallen.
\end{werkwijze}

\begin{oefeningen}
\oef{h1-combinaties}
\end{oefeningen}

%========================================================================================
\section{Relatieve frequentie}

\subsection{Een kans schatten}

Is een punaise op een tafel geworpen, dan landt ze op haar kop of op haar punt. Die twee uitkomsten zijn niet even waarschijnlijk, en een formule voor de kans is er niet. Dan doe je het experiment vaak, en kijk je hoe vaak elke uitkomst voorkomt.

\begin{theorie}
De \textbf{\textcolor{accent}{relatieve frequentie}} van een gebeurtenis $A$ na $n$ herhalingen is \quad $\formule{\dfrac{\text{aantal keer } A}{n}}$
\end{theorie}

Een relatieve frequentie ligt altijd tussen 0 en 1, en de relatieve frequenties van alle uitkomsten zijn samen 1, net als kansen. Je schrijft ze ook als percentage.

\voorbeeld{}{Een punaise wordt 1000 keer geworpen en landt 412 keer op haar punt. Schat de kans op \textit{punt}. Hoe vaak verwacht je \textit{punt} bij 250 worpen?\label{vb:h1-frequentie}}{%
\uitleg{De relatieve frequentie na 1000 worpen is een schatting van de kans. Met die kans schatten we het aantal bij een nieuwe reeks.}
\begin{stappen}
P(\text{punt}) \approx \frac{412}{1000} = 0{,}412 & De relatieve frequentie. \\[2.5ex]
250 \cdot 0{,}412 = 103 & Ongeveer 103 keer, maar het werkelijke aantal schommelt.
\end{stappen}}

\begin{oefeningen}
\oef{h1-frequentie}

\oef{h1-frequentie-2}
\end{oefeningen}

\needspace{16\baselineskip}
\subsection{De wet van de grote aantallen}

\naast{\begin{minipage}{9cm}
Hoe goed is zo'n schatting? Bij een dobbelsteen kennen we de kans op een zes, en kunnen we het nagaan. De figuur toont een simulatie: een computer gooide 2000 keer met een dobbelsteen, en na elke worp is de relatieve frequentie van een zes berekend. In het begin schommelt ze sterk: na 20 worpen is ze $0{,}25$. Hoe meer worpen, hoe dichter ze bij $\frac16 \approx 0{,}167$ blijft. Eén worp meer of minder verandert na 2000 worpen bijna niets meer aan de verhouding.
\end{minipage}}{\begin{tikzpicture}[x=.0026cm,y=14cm,font=\scriptsize]
  \draw[black!14,very thin] (0,0) grid[xstep=500,ystep=.05] (2000,.3);
  \draw[->,grijs] (0,0) -- (2120,0) node[right,text=black] {$n$};
  \draw[->,grijs] (0,0) -- (0,.32);
  \foreach \t in {500,1000,1500,2000}{\node[below] at (\t,0) {\t};}
  \foreach \t/\l in {.1/{0,1},.2/{0,2},.3/{0,3}}{\node[left] at (0,\t) {\l};}
  \draw[accent,dashed] (0,1/6) -- (2000,1/6) node[right,text=accent] {$\tfrac16$};
  \draw[fel!80!black,thick] plot[smooth] coordinates {(10,.1) (20,.25) (40,.15) (60,.183) (100,.23) (150,.18) (200,.16) (300,.157) (400,.165) (600,.167) (800,.166) (1000,.168) (1200,.168) (1400,.164) (1600,.165) (1800,.166) (2000,.168)};
\end{tikzpicture}}

\begin{theorie}
\begin{stelling}{Wet van de grote aantallen}
Herhaal je een kansexperiment heel vaak, onafhankelijk en onder dezelfde omstandigheden, dan nadert de relatieve frequentie van $A$ tot $\formule{P(A)}$.
\end{stelling}
\end{theorie}

Het bewijs, en wat \textit{nadert} hier precies betekent, valt buiten deze cursus. De wet zegt wel waarom je een kans die je niet kunt berekenen, mag schatten met een relatieve frequentie na veel herhalingen. Ze zegt ook iets niet: ze gaat over de verhouding op lange termijn, niet over de volgende worp.

\voorbeeld{}{Een eerlijke munt landde vijf keer na elkaar op kop. Is munt bij de zesde worp nu waarschijnlijker dan kop?\label{vb:h1-gokker}}{%
\uitleg{Veel mensen denken dat munt \textit{aan de beurt is}, zodat de verhouding weer klopt. Kijk wat er echt gebeurt in de volgende 1000 worpen.}
\begin{stappenbreed}
P(\text{munt bij worp 6}) = \frac12 & Een munt heeft geen geheugen: de vorige worpen veranderen niets aan deze. \\[1ex]
\frac{5 + 500}{5 + 1000} \approx 0{,}502 & Na nog 1000 worpen verwacht je ongeveer 500 keer kop. De relatieve frequentie van kop is dan bijna $\frac12$.
\end{stappenbreed}
De vijf extra keer kop worden dus niet \textit{ingehaald} door extra munt: ze wegen gewoon steeds minder door in een steeds grotere reeks.}

\begin{oefeningen}
\oef{h1-gokker}
\end{oefeningen}

%========================================================================================
\section{Rekenregels}

Zowel met Laplace als met relatieve frequenties is een kans een aantal gedeeld door een totaal. Daaruit volgen een paar regels die altijd gelden. We bewijzen ze telkens voor de wet van Laplace. Voor relatieve frequenties is het bewijs hetzelfde, met $n$ in de plaats van $\#\Omega$.

\subsection{Grenzen van een kans}

\begin{theorie}
\begin{stelling}{Grenzen}
Voor elke gebeurtenis $A$ is \quad $\formule{0 \leqslant P(A) \leqslant 1}$, \quad met \quad $\formule{P(\emptyset) = 0}$ \quad en \quad $\formule{P(\Omega) = 1}$
\end{stelling}

\begin{bewijs}
$A$ is een deel van $\Omega$, dus $0 \leqslant \#A \leqslant \#\Omega$. Deel door $\#\Omega > 0$:
\bewijsregel{0 \leqslant \frac{\#A}{\#\Omega} \leqslant 1 .}
De onmogelijke gebeurtenis heeft geen uitkomsten, dus $P(\emptyset) = \frac{0}{\#\Omega} = 0$, en $P(\Omega) = \frac{\#\Omega}{\#\Omega} = 1$.
\end{bewijs}
\end{theorie}

Een kans van 0 betekent \textit{onmogelijk}, een kans van 1 \textit{zeker}. Vind je een kans groter dan 1 of kleiner dan 0, dan heb je een rekenfout gemaakt. Die controle is snel, en ze vangt een klassieke fout op.

\voorbeeld{}{Sara gooit drie keer met een munt. Ze zegt: \textit{elke worp geeft kop met kans $\frac12$, dus de kans op minstens één keer kop is $\frac12 + \frac12 + \frac12 = \frac32$.} Wat is er fout, en wat is de juiste kans?\label{vb:h1-grenzen}}{%
\uitleg{Een kans van $\frac32$ kan niet, dus de redenering klopt niet. We tellen de rijtjes zelf.}
\begin{stappenbreed}
\#\Omega = 2 \cdot 2 \cdot 2 = 8 & De rijtjes $KKK, KKM, \ldots, MMM$ zijn even waarschijnlijk. \\[1ex]
P(\text{minstens één kop}) = \frac78 & Alleen $MMM$ bevat geen kop.
\end{stappenbreed}
Sara telt een rijtje als $KKM$ meerdere keren: één keer bij \textit{kop bij worp 1} en één keer bij \textit{kop bij worp 2}. Die gebeurtenissen zijn niet onverenigbaar, en dan mag je hun kansen niet zomaar optellen.}

\begin{oefeningen}
\oef{h1-grenzen}
\end{oefeningen}

\needspace{16\baselineskip}
\subsection{De somregel voor onverenigbare gebeurtenissen}

Wanneer mag je kansen wel optellen? Als de gebeurtenissen niet samen kunnen optreden: dan telt geen enkele uitkomst dubbel.

\begin{theorie}
\begin{stelling}{Somregel voor onverenigbare gebeurtenissen}
Zijn $A$ en $B$ onverenigbaar, dan is \quad $\formule{P(A \cup B) = P(A) + P(B)}$
\end{stelling}

\begin{bewijs}
Omdat $A \cap B = \emptyset$, ligt elke uitkomst van $A \cup B$ in precies één van beide. Tel je de uitkomsten van $A$ en die van $B$, dan tel je elke uitkomst van $A \cup B$ precies één keer:
\bewijsregel{\#(A \cup B) = \#A + \#B .}
Deel door $\#\Omega$: dan staat er $P(A \cup B) = P(A) + P(B)$.
\end{bewijs}
\end{theorie}

Kortom: \textit{of} wordt \textit{plus}, als de gebeurtenissen elkaar uitsluiten. Dat geldt ook voor meer dan twee gebeurtenissen die twee aan twee onverenigbaar zijn: tel alle kansen op.

\voorbeeld{}{Je gooit met twee dobbelstenen. Bij een som van 2, 3 of 12 verlies je meteen. Wat is de kans daarop?\label{vb:h1-somregel}}{%
\uitleg{Een worp heeft maar één som, dus de drie gebeurtenissen zijn twee aan twee onverenigbaar. We tellen hun kansen op.}
\begin{stappenbreed}
P(\text{som } 2) = \frac{1}{36}, \ P(\text{som } 3) = \frac{2}{36}, \ P(\text{som } 12) = \frac{1}{36} & Uit het rooster: $(1, 1)$; $(1, 2), (2, 1)$; $(6, 6)$. \\[2ex]
P = \frac{1 + 2 + 1}{36} = \frac{4}{36} = \frac19 & De somregel, drie keer.
\end{stappenbreed}}

\begin{oefeningen}
\oef{h1-onverenigbaar}
\end{oefeningen}

\needspace{16\baselineskip}
\subsection{De complementregel}

Een gebeurtenis en haar complement zijn onverenigbaar, en samen vormen ze $\Omega$. Met de somregel volgt daaruit een regel die veel werk spaart.

\begin{theorie}
\begin{stelling}{Complementregel}
Voor elke gebeurtenis $A$ is \quad $\formule{P(\compl{A}) = 1 - P(A)}$
\end{stelling}

\begin{bewijs}
$A$ en $\compl{A}$ zijn onverenigbaar en $A \cup \compl{A} = \Omega$. De somregel en $P(\Omega) = 1$ geven
\bewijsregel{P(A) + P(\compl{A}) = P(\Omega) = 1 ,}
waaruit $P(\compl{A}) = 1 - P(A)$.
\end{bewijs}
\end{theorie}

De regel is vooral handig bij \textit{minstens één}. Die gebeurtenis omvat veel gevallen: precies één, precies twee, \ldots{} Het tegendeel is \textit{geen enkele}, en dat is vaak één eenvoudig geval.

\voorbeeld{}{Je gooit vier keer met een dobbelsteen. Wat is de kans op minstens één zes?\label{vb:h1-complement}}{%
\uitleg{\textit{Minstens één zes} omvat veel gevallen. Het tegendeel is \textit{geen enkele zes}: bij elke worp een van de 5 andere getallen.}
\begin{stappen}
P(\text{geen zes}) = \frac{5^4}{6^4} = \frac{625}{1296} & Laplace, met $6^4$ even waarschijnlijke rijtjes, waarvan $5^4$ zonder zes. \\[3ex]
P(\text{minstens één zes}) = 1 - \frac{625}{1296} \approx 0{,}518 & Het complement: iets meer dan een half.
\end{stappen}}

Wie hier $4 \cdot \frac16 = \frac23$ rekent, maakt de fout van Sara: \textit{een zes bij worp 1} en \textit{een zes bij worp 2} zijn niet onverenigbaar. Bij zeven worpen zou die redenering zelfs $\frac76$ geven. Let ook op het tegendeel: dat van \textit{minstens één} is \textit{geen enkele}, niet \textit{precies één}.

\begin{werkwijze}
Bij \textit{minstens één}: bereken de kans op het tegendeel, \textit{geen enkele}, en trek die af van 1.
\end{werkwijze}

\begin{oefeningen}
\oef{h1-minstens}
\end{oefeningen}

\needspace{16\baselineskip}
\subsection{De algemene somregel}

Zijn $A$ en $B$ niet onverenigbaar, dan telt $P(A) + P(B)$ de uitkomsten in de overlap twee keer. Bij een dobbelsteen is de kans op \textit{even of minstens 4} niet $\frac12 + \frac12 = 1$: dan zou het zeker zijn, terwijl een 1 of een 3 kan. De overlap moet er één keer af.

\begin{theorie}
\begin{stelling}{Algemene somregel}
\begin{minipage}[t]{.64\linewidth}
\vspace{0pt}
Voor alle gebeurtenissen $A$ en $B$ is

\vspace{1mm}
\qquad $\formule{P(A \cup B) = P(A) + P(B) - P(A \cap B)}$
\end{minipage}\hfill
\begin{minipage}[t]{.3\linewidth}
\vspace{0pt}
\begin{tikzpicture}[font=\scriptsize,scale=.62]
  \fill[accent!45] (1.15,0) circle (.85); \fill[accent!12] (0,0) circle (.85);
  \draw[accent,thick] (0,0) circle (.85); \draw[accent,thick] (1.15,0) circle (.85);
  \node[text=accent] at (-.35,0) {$A$}; \node[text=accent] at (1.6,0) {$B$};
  \node[below] at (1.4,-.85) {$B \cap \compl{A}$};
\end{tikzpicture}
\end{minipage}
\end{stelling}

\begin{bewijs}
$A \cup B$ valt uiteen in $A$ en het deel van $B$ buiten $A$, dus $B \cap \compl{A}$. Ook $B$ zelf valt uiteen in $A \cap B$ en $B \cap \compl{A}$. Telkens zijn de delen onverenigbaar, dus de somregel geeft
\bewijsregel{P(A \cup B) = P(A) + P(B \cap \compl{A}) \qquad\text{en}\qquad P(B) = P(A \cap B) + P(B \cap \compl{A}) .}
Uit de tweede gelijkheid is $P(B \cap \compl{A}) = P(B) - P(A \cap B)$. Vul dat in in de eerste.
\end{bewijs}
\end{theorie}

Zijn $A$ en $B$ onverenigbaar, dan is $P(A \cap B) = 0$ en vind je de somregel van daarnet terug. In het voorbeeld van de dobbelsteen is $P(\text{even of minstens 4}) = \frac36 + \frac36 - \frac26 = \frac46$: de uitkomsten 4 en 6 zijn één keer te veel geteld.

\voorbeeld{}{Je trekt een kaart uit een spel van 52. Wat is de kans op een harten of een aas?\label{vb:h1-som}}{%
\uitleg{De gebeurtenissen overlappen: de hartenaas is beide. We gebruiken de algemene somregel.}
\begin{stappen}
P(H \cup A) = \frac{13}{52} + \frac{4}{52} - \frac{1}{52} = \frac{16}{52} = \frac{4}{13} & De hartenaas zou anders twee keer geteld zijn.
\end{stappen}}

\begin{werkwijze}
Bij \textit{$A$ of $B$}: ga na of $A$ en $B$ samen kunnen optreden. Zo niet, tel de kansen op. Zo wel, tel ze op en trek de kans op \textit{$A$ en $B$} af. Controleer dat de uitkomst tussen 0 en 1 ligt.
\end{werkwijze}

\begin{oefeningen}
\oef{h1-somregel}

\oef{h1-regels}
\end{oefeningen}

%========================================================================================
\opdrachtenkop

\begin{oefeningen}
\needspace{12\baselineskip}
\oef{h1-o13}

\oef{h1-o1}

\oef{h1-o16}

\oef{h1-o3}

\oef{h1-o20}

\oef{h1-o2}

\needspace{14\baselineskip}
\oef{h1-o14}

\oef{h1-o15}

\oef{h1-o17}

\oef{h1-o4}

\oef{h1-o5}

\needspace{19\baselineskip}
\oef{h1-o19}

\oef{h1-o23}

\oef{h1-o7}

\needspace{10\baselineskip}
\oef{h1-o18}

\oef{h1-o22}

\oef{h1-o24}

\oef{h1-o6}

\needspace{10\baselineskip}
\oef{h1-o21}

\oef{h1-o8}

\oef{h1-o9}

\oef{h1-o10}

\oef{h1-o26}

\oef{h1-o11}

\oef{h1-o25}

\oef{h1-o27}

\oef{h1-o12}
\end{oefeningen}

%----------------------------------------------------------------------------------------
% KOPPEL: verbanden met andere leerstof
%----------------------------------------------------------------------------------------
\needspace{20\baselineskip}
\opdrachtsectie{Koppel}

\lkr{De eerste oefening gebruikt de oppervlakte van een cirkel. De tweede gebruikt variaties uit de cursus Telproblemen (zelfde jaar): kijk na of die al gezien is.}

\begin{oefeningen}
\oef{h1-koppel-1}

\oef{h1-koppel-2}
\end{oefeningen}

%----------------------------------------------------------------------------------------
\opdrachtsectie{Test jezelf}

\begin{oefeningen}
\oef{h1-t1}

\oef{h1-t2}

\oef{h1-t3}

\oef{h1-t4}

\oef{h1-t5}

\oef{h1-t6}

\oef{h1-t7}
\end{oefeningen}

\herhaalschema{%
  \fv{oef:h1-t1}{vb:h1-laplace}{oef:h1-laplace}%
  \fv{oef:h1-t2}{vb:h1-complement}{oef:h1-minstens}%
  \fv{oef:h1-t3}{vb:h1-som}{oef:h1-regels}%
  \fv{oef:h1-t4}{vb:h1-kaarten}{oef:h1-combinaties}%
  \fv{oef:h1-t5}{vb:h1-frequentie}{oef:h1-frequentie}%
  \fv{oef:h1-t6}{vb:h1-gebeurtenis}{oef:h1-gebeurtenis}%
  \fv{oef:h1-t7}{vb:h1-som}{oef:h1-somregel}}

\ifzelfstudie
\section*{Maak je eigen samenvatting}
Sluit dit boek en schrijf in je schrift: wat een kansexperiment, een uitkomstenverzameling en een gebeurtenis zijn, de notaties $\compl{A}$, $A \cap B$ en $A \cup B$, de wet van Laplace met haar voorwaarde, de relatieve frequentie en de wet van de grote aantallen, en de rekenregels voor kansen. Vergelijk daarna met de kaders.
\fi
